\section{De la resolución de problemas a la investigación matemática}

AI for Math puede dividirse en distintos niveles: resolver problemas, demostración formal, completado de demostraciones, generación de conjeturas y colaboración en investigación. Cuanto más se avanza, más se acerca a la investigación matemática real y más difícil es evaluar. Los problemas de concurso tienen respuesta, las demostraciones en Lean tienen un verifier, mientras que en los problemas de investigación hay que juzgar «si vale la pena abordarlos».

\begin{figure}[H]
\centering
\includegraphics[width=0.96\textwidth]{ai-for-math-roadmap.png}
\caption{Hoja de ruta de AI for Math: de los problemas matemáticos a la investigación matemática.}
\end{figure}

\begin{knowledgebox}{Lectura de la figura: por qué es más difícil cuanto más a la derecha}
La resolución de problemas suele tener una respuesta; la demostración formal se verifica con un proof assistant; el completado de demostraciones tiene objetivos locales; en cambio, la generación de conjeturas y la colaboración en investigación requieren sentido estético, criterio para elegir direcciones y una estructura teórica de largo plazo. Cuanto más cerca de la investigación, menos se puede depender solo de un benchmark.
\end{knowledgebox}

\subsection{La intuición del matemático}

La segunda mitad de la entrevista trata sobre «la intuición del matemático». La intuición no es algo místico, sino un juicio condensado sobre estructuras, contraejemplos, rutas de demostración y estética, formado tras un entrenamiento prolongado. La IA puede aprender patrones a partir de grandes cantidades de datos y de retroalimentación, pero que llegue a formar algo parecido a la intuición de un matemático depende de si logra integrar la retroalimentación formal, el razonamiento informal y los objetivos de investigación.

\begin{knowledgebox}{La intuición puede entrenarse, pero no equivale a ajustar datos}
La intuición matemática surge de la elección de problemas, la experiencia con los fracasos, las analogías estructurales y el juicio estético. La IA puede aprender de bibliotecas de demostraciones, artículos, retroalimentación interactiva y correcciones de matemáticos, pero necesita un ciclo cerrado de largo plazo, no datos supervisados de una sola vez.
\end{knowledgebox}

\subsection{Un viaje a la Luna o tiene éxito o fracasa}

La expresión de la entrevista «un viaje a la Luna o tiene éxito o fracasa» refleja el carácter de alto riesgo del emprendimiento en AI for Math. La investigación matemática no es como una aplicación común, en la que se pueden validar indicadores comerciales en pasos pequeños: o impulsa de verdad la capacidad de demostrar y descubrir, o es solo una demo vistosa. La dificultad de un Neo Lab como Axiom consiste en que debe ocuparse a la vez de la ciencia, el producto, el equipo y el financiamiento.

\begin{warningbox}{La evaluación de AI for Math no puede basarse solo en el financiamiento}
Las grandes rondas de financiamiento y la incorporación de matemáticos de primer nivel son señales, pero no resultados. Los resultados reales dependen de si se producen demostraciones verificables, se acelera el trabajo de los matemáticos y se forman bibliotecas de teoremas reutilizables y modelos de colaboración en investigación.
\end{warningbox}

\subsection{Resumen de la sección}

La ruta de AI for Math va de la resolución de problemas a la colaboración en investigación. Cuanto más se acerca a la investigación, más necesita la intuición del matemático, sistemas formales y retroalimentación de largo plazo, y no un único benchmark.
